% GENERATED FILE. Edit canonical lessons, then run npm run generate:books.
% canonical-source-hash: fnv1a64:2f3cde9ab88ce6e0
% canonical-lessons: FR-C139-rien, FR-C139-aucun, FR-C139-non-plus, FR-C139-lequel, FR-C139-nimporte-quoi

\chapter{Nothing, None, Neither, Which One, Anything at All}
\label{ch:fr-nothing-none-neither-which-one-anything-at-all}
\hlchaptermodality{139}

\begin{chapteropening}
\textbf{By the end of this chapter:} \emph{I can say that there is nothing and none, ask which one, and answer \textquotedblleft{}me neither\textquotedblright{}.}

Complete the last lesson of chapter 139: I can say that there is nothing and none, ask which one, and answer \textquotedblleft{}me neither\textquotedblright{}.
\end{chapteropening}

\section[rien]{rien — nothing}

\label{lesson:FR-C139-rien}



\begin{quote}
Before the new one: say the French for to drive, then the French for to get dressed.
\end{quote}

\subsection*{You'll want to know: rien}
\textbf{rien} — \textquotedblleft{}nothing\textquotedblright{}. With a verb it takes \textbf{ne} in front: \textbf{Je ne comprends rien.} — \textquotedblleft{}I understand nothing.\textquotedblright{} On its own it answers a question: \textbf{Rien.}

\begin{grammarlens}[title={What you've built}]
One more word for this chapter.
\end{grammarlens}

\subsection*{Your turn}
\begin{itemize}
\raggedright
  \item \emph{Say it:} \emph{rien}
  \item \emph{Say it:} \emph{rien}, once more
  \item \emph{Say it:} say \emph{rien}
  \item \emph{Recall:} say the French for to drive, then the French for to get dressed, then say \emph{rien} again
\end{itemize}

\begin{tcolorbox}[breakable,colback=teal!4,colframe=teal!35!black,title={Before you move on}]
What does rien mean? (\textquotedblleft{}Nothing\textquotedblright{}.) Say it once more.
\end{tcolorbox}
\section[aucun / aucune]{aucun / aucune — none, not a single}

\label{lesson:FR-C139-aucun}



\begin{quote}
Before the new one: say the French for to get dressed, then the French for nothing.
\end{quote}

\subsection*{You'll want to know: aucun / aucune}
\textbf{aucun / aucune} — \textquotedblleft{}none, not a single\textquotedblright{}. It agrees with its noun and also takes \textbf{ne}: \textbf{Il n'y a aucun train lundi.} — \textquotedblleft{}There is no train at all on Monday.\textquotedblright{}

\begin{grammarlens}[title={What you've built}]
One more word for this chapter.
\end{grammarlens}

\subsection*{Your turn}
\begin{itemize}
\raggedright
  \item \emph{Say it:} \emph{aucun / aucune}
  \item \emph{Say it:} \emph{aucun / aucune}, once more
  \item \emph{Say it:} say \emph{aucun / aucune}
  \item \emph{Recall:} say the French for to get dressed, then the French for nothing, then say \emph{aucun / aucune} again
\end{itemize}

\begin{tcolorbox}[breakable,colback=teal!4,colframe=teal!35!black,title={Before you move on}]
What does aucun / aucune mean? (\textquotedblleft{}None, not a single\textquotedblright{}.) Say it once more.
\end{tcolorbox}
\section[non plus]{non plus — neither, not either}

\label{lesson:FR-C139-non-plus}



\begin{quote}
Before the new one: say the French for nothing, then the French for none.
\end{quote}

\subsection*{You'll want to know: non plus}
\textbf{non plus} — \textquotedblleft{}neither, not either\textquotedblright{}. The answer to a negative: \textbf{Je ne travaille pas samedi. — Moi non plus.} — \textquotedblleft{}Me neither.\textquotedblright{} After a yes, French says \textbf{moi aussi} instead.

\begin{grammarlens}[title={What you've built}]
One more word for this chapter.
\end{grammarlens}

\subsection*{Your turn}
\begin{itemize}
\raggedright
  \item \emph{Say it:} \emph{non plus}
  \item \emph{Say it:} \emph{non plus}, once more
  \item \emph{Say it:} say \emph{non plus}
  \item \emph{Recall:} say the French for nothing, then the French for none, then say \emph{non plus} again
\end{itemize}

\begin{tcolorbox}[breakable,colback=teal!4,colframe=teal!35!black,title={Before you move on}]
What does non plus mean? (\textquotedblleft{}Neither, not either\textquotedblright{}.) Say it once more.
\end{tcolorbox}
\section[lequel / laquelle]{lequel / laquelle — which one}

\label{lesson:FR-C139-lequel}



\begin{quote}
Before the new one: say the French for none, then the French for neither.
\end{quote}

\subsection*{You'll want to know: lequel / laquelle}
\textbf{lequel / laquelle} — \textquotedblleft{}which one\textquotedblright{}. A question that picks from a set, and it agrees with what is picked: \textbf{Lequel ?} for a train, \textbf{Laquelle ?} for a chambre.

\begin{grammarlens}[title={What you've built}]
One more word for this chapter.
\end{grammarlens}

\subsection*{Your turn}
\begin{itemize}
\raggedright
  \item \emph{Say it:} \emph{lequel / laquelle}
  \item \emph{Say it:} \emph{lequel / laquelle}, once more
  \item \emph{Say it:} say \emph{lequel / laquelle}
  \item \emph{Recall:} say the French for none, then the French for neither, then say \emph{lequel / laquelle} again
\end{itemize}

\begin{tcolorbox}[breakable,colback=teal!4,colframe=teal!35!black,title={Before you move on}]
What does lequel / laquelle mean? (\textquotedblleft{}Which one\textquotedblright{}.) Say it once more.
\end{tcolorbox}
\section[n'importe quoi]{n'importe quoi — anything at all}

\label{lesson:FR-C139-nimporte-quoi}



\begin{quote}
Before the new one: say the French for neither, then the French for which one.
\end{quote}

\subsection*{You'll want to know: n'importe quoi}
\textbf{n'importe quoi} — \textquotedblleft{}anything at all\textquotedblright{}. \textbf{Je lis n'importe quoi.} — \textquotedblleft{}I read anything at all.\textquotedblright{} Said of what someone claims, it means \textquotedblleft{}nonsense\textquotedblright{}.

\begin{grammarlens}[title={What you've built}]
One more word for this chapter.
\end{grammarlens}

\subsection*{Your turn}
\begin{itemize}
\raggedright
  \item \emph{Say it:} \emph{n'importe quoi}
  \item \emph{Say it:} \emph{n'importe quoi}, once more
  \item \emph{Say it:} say \emph{n'importe quoi}
  \item \emph{Recall:} say the French for neither, then the French for which one, then say \emph{n'importe quoi} again
\end{itemize}

\begin{tcolorbox}[breakable,colback=teal!4,colframe=teal!35!black,title={Before you move on}]
What does n'importe quoi mean? (\textquotedblleft{}Anything at all\textquotedblright{}.) Say it once more.
\end{tcolorbox}
